% GENERATED FILE. Edit canonical lessons, then run npm run generate:books.
% canonical-source-hash: fnv1a64:30ea5bbcebd85b98
% canonical-lessons: TE-C492-porabadu, TE-C492-tadabadu, TE-C492-ancana-veyu, TE-C492-tappatadugu, TE-C492-gandaragolam

\chapter{Mistakes and Mix-ups}
\label{ch:te-mistakes-and-mix-ups}
\hlchaptermodality{492}

\begin{chapteropening}
\textbf{By the end of this chapter:} \emph{I can say to be mistaken, to falter, to estimate, a false step and confusion.}

Complete the last lesson of chapter 492: I can say to be mistaken, to falter, to estimate, a false step and confusion.
\end{chapteropening}

\section[\texorpdfstring{\te{పొరబడు} (porabaḍu)}{పొరబడు (porabaḍu)}]{\te{పొరబడు} (porabaḍu) — to be mistaken}

\label{lesson:TE-C492-porabadu}



\begin{quote}
Before the new one: say the Telugu for to crack, then the Telugu for to deteriorate.
\end{quote}

\subsection*{You'll want to know: \te{పొరబడు}}
\textbf{\te{పొరబడు}} — \emph{porabaḍu} — \textquotedblleft{}to be mistaken\textquotedblright{}.

\begin{quote}
\te{ఆ} \te{బస్సు} \te{మనదే} \te{అని} \te{పొరబడ్డాను}.
\end{quote}

\begin{quote}
\emph{ā bassu manadē ani porabaḍḍānu.}
\end{quote}

\begin{quote}
I mistook that bus for ours.
\end{quote}

Past (I …): \te{పొరబడ్డాను}, \emph{porabaḍḍānu}.

\te{పొర}, an error, + \te{పడు}. \te{పొరపాటు}, a mistake, comes from the same \te{పొర}.

\begin{cousinweb}
\textbf{\te{పొరబడు}} (\emph{porabaḍu}) is native Dravidian: DEDR 4293 lists Telugu \emph{pora}, \textquotedblleft{}error, misunderstanding\textquotedblright{}, and \emph{porabaḍu}, \textquotedblleft{}to go the wrong way\textquotedblright{}.
\end{cousinweb}

\begin{grammarlens}[title={What you've built}]
One more word for getting around.
\end{grammarlens}

\subsection*{Your turn}
\begin{itemize}
\raggedright
  \item \emph{Say it:} \emph{porabaḍu}
  \item \emph{Say it:} \emph{porabaḍu}, once more
  \item \emph{Say it:} \emph{ā bassu manadē ani porabaḍḍānu.} — I mistook that bus for ours.
  \item \emph{Recall:} say the Telugu for to crack, then the Telugu for to deteriorate, then say \emph{porabaḍu} again
\end{itemize}

\begin{tcolorbox}[breakable,colback=teal!4,colframe=teal!35!black,title={Before you move on}]
What does \te{పొరబడు} mean? (\textquotedblleft{}To be mistaken\textquotedblright{}.) Say it once more.
\end{tcolorbox}
\section[\texorpdfstring{\te{తడబడు} (taḍabaḍu)}{తడబడు (taḍabaḍu)}]{\te{తడబడు} (taḍabaḍu) — to falter; to stumble}

\label{lesson:TE-C492-tadabadu}



\begin{quote}
Before the new one: say the Telugu for to deteriorate, then the Telugu for to be mistaken.
\end{quote}

\subsection*{You'll want to know: \te{తడబడు}}
\textbf{\te{తడబడు}} — \emph{taḍabaḍu} — \textquotedblleft{}to falter; to stumble\textquotedblright{}.

\begin{quote}
\te{కొత్త} \te{భాషలో} \te{మాట్లాడుతూ} \te{తడబడ్డాను}.
\end{quote}

\begin{quote}
\emph{kotta bhāṣalō māṭlāḍutū taḍabaḍḍānu.}
\end{quote}

\begin{quote}
I faltered while speaking the new language.
\end{quote}

Past (I …): \te{తడబడ్డాను}, \emph{taḍabaḍḍānu}.

Said of the feet, the tongue and a nervous speaker.

\begin{cousinweb}
\textbf{\te{తడబడు}} (\emph{taḍabaḍu}) is native Dravidian: DEDR 3033 lists Telugu \emph{taḍabaḍu}, \textquotedblleft{}to be confused, bewildered\textquotedblright{}, with Tamil \emph{taṭavu}, \textquotedblleft{}to be halting\textquotedblright{}.
\end{cousinweb}

\begin{grammarlens}[title={What you've built}]
One more word for getting around.
\end{grammarlens}

\subsection*{Your turn}
\begin{itemize}
\raggedright
  \item \emph{Say it:} \emph{taḍabaḍu}
  \item \emph{Say it:} \emph{taḍabaḍu}, once more
  \item \emph{Say it:} \emph{kotta bhāṣalō māṭlāḍutū taḍabaḍḍānu.} — I faltered while speaking the new language.
  \item \emph{Recall:} say the Telugu for to deteriorate, then the Telugu for to be mistaken, then say \emph{taḍabaḍu} again
\end{itemize}

\begin{tcolorbox}[breakable,colback=teal!4,colframe=teal!35!black,title={Before you move on}]
What does \te{తడబడు} mean? (\textquotedblleft{}To falter; to stumble\textquotedblright{}.) Say it once more.
\end{tcolorbox}
\section[\texorpdfstring{\te{అంచనా} \te{వేయు} (añcanā vēyu)}{అంచనా వేయు (añcanā vēyu)}]{\te{అంచనా} \te{వేయు} (añcanā vēyu) — to estimate}

\label{lesson:TE-C492-ancana-veyu}



\begin{quote}
Before the new one: say the Telugu for to be mistaken, then the Telugu for to falter.
\end{quote}

\subsection*{You'll want to know: \te{అంచనా} \te{వేయు}}
\textbf{\te{అంచనా} \te{వేయు}} — \emph{añcanā vēyu} — \textquotedblleft{}to estimate\textquotedblright{}.

\begin{quote}
\te{ప్రయాణానికి} \te{ఎంత} \te{ఖర్చు} \te{అవుతుందో} \te{అంచనా} \te{వేశాం}.
\end{quote}

\begin{quote}
\emph{prayāṇāniki enta kharcu avutundō añcanā vēśāṁ.}
\end{quote}

\begin{quote}
We estimated how much the trip would cost.
\end{quote}

Past (I …): \te{అంచనా} \te{వేశాను}, \emph{añcanā vēśānu}.

\te{అంచనా}, an estimate, + \te{వేయు}, to put. The weather news also gives \te{అంచనాలు}, forecasts.

\begin{cousinweb}
\textbf{\te{అంచనా} \te{వేయు}} (\emph{añcanā vēyu}) is a Telugu phrase: \te{అంచనా}, \textquotedblleft{}an estimate\textquotedblright{}, whose origin is not settled, and the native \te{వేయు}, \textquotedblleft{}to put\textquotedblright{} (DEDR 5549).
\end{cousinweb}

\begin{grammarlens}[title={What you've built}]
One more word for getting around.
\end{grammarlens}

\subsection*{Your turn}
\begin{itemize}
\raggedright
  \item \emph{Say it:} \emph{añcanā vēyu}
  \item \emph{Say it:} \emph{añcanā vēyu}, once more
  \item \emph{Say it:} \emph{prayāṇāniki enta kharcu avutundō añcanā vēśāṁ.} — We estimated how much the trip would cost.
  \item \emph{Recall:} say the Telugu for to be mistaken, then the Telugu for to falter, then say \emph{añcanā vēyu} again
\end{itemize}

\begin{tcolorbox}[breakable,colback=teal!4,colframe=teal!35!black,title={Before you move on}]
What does \te{అంచనా} \te{వేయు} mean? (\textquotedblleft{}To estimate\textquotedblright{}.) Say it once more.
\end{tcolorbox}
\section[\texorpdfstring{\te{తప్పటడుగు} (tappaṭaḍugu)}{తప్పటడుగు (tappaṭaḍugu)}]{\te{తప్పటడుగు} (tappaṭaḍugu) — a false step; a blunder}

\label{lesson:TE-C492-tappatadugu}



\begin{quote}
Before the new one: say the Telugu for to falter, then the Telugu for to estimate.
\end{quote}

\subsection*{You'll want to know: \te{తప్పటడుగు}}
\textbf{\te{తప్పటడుగు}} — \emph{tappaṭaḍugu} — \textquotedblleft{}a false step; a blunder\textquotedblright{}.

\begin{quote}
\te{తొందరపడి} \te{ఆ} \te{ఉద్యోగం} \te{వదలడం} \te{నా} \te{తప్పటడుగు}.
\end{quote}

\begin{quote}
\emph{tondarapaḍi ā udyōgaṁ vadalaḍaṁ nā tappaṭaḍugu.}
\end{quote}

\begin{quote}
Leaving that job in a hurry was my false step.
\end{quote}

\te{తప్పు}, wrong, + \te{అడుగు}, a step. A toddler's first wobbly steps are \te{తప్పటడుగులు} too.

\begin{cousinweb}
\textbf{\te{తప్పటడుగు}} (\emph{tappaṭaḍugu}) is native Dravidian: \te{తప్పు}, \textquotedblleft{}to slip, to go wrong\textquotedblright{} (DEDR 3071), and \te{అడుగు}, \textquotedblleft{}a footstep, a step\textquotedblright{} (DEDR 72), joined by -\te{ట}-.
\end{cousinweb}

\begin{grammarlens}[title={What you've built}]
One more word for getting around.
\end{grammarlens}

\subsection*{Your turn}
\begin{itemize}
\raggedright
  \item \emph{Say it:} \emph{tappaṭaḍugu}
  \item \emph{Say it:} \emph{tappaṭaḍugu}, once more
  \item \emph{Say it:} \emph{tondarapaḍi ā udyōgaṁ vadalaḍaṁ nā tappaṭaḍugu.} — Leaving that job in a hurry was my false step.
  \item \emph{Recall:} say the Telugu for to falter, then the Telugu for to estimate, then say \emph{tappaṭaḍugu} again
\end{itemize}

\begin{tcolorbox}[breakable,colback=teal!4,colframe=teal!35!black,title={Before you move on}]
What does \te{తప్పటడుగు} mean? (\textquotedblleft{}A false step; a blunder\textquotedblright{}.) Say it once more.
\end{tcolorbox}
\section[\texorpdfstring{\te{గందరగోళం} (gandaragōḷaṁ)}{గందరగోళం (gandaragōḷaṁ)}]{\te{గందరగోళం} (gandaragōḷaṁ) — confusion; a muddle}

\label{lesson:TE-C492-gandaragolam}



\begin{quote}
Before the new one: say the Telugu for to estimate, then the Telugu for a false step.
\end{quote}

\subsection*{You'll want to know: \te{గందరగోళం}}
\textbf{\te{గందరగోళం}} — \emph{gandaragōḷaṁ} — \textquotedblleft{}confusion; a muddle\textquotedblright{}.

\begin{quote}
\te{టిక్కెట్లు} \te{తారుమారు} \te{కావడంతో} \te{స్టేషన్లో} \te{గందరగోళం} \te{మొదలైంది}.
\end{quote}

\begin{quote}
\emph{ṭikkeṭlu tārumāru kāvaḍantō sṭēṣanlō gandaragōḷaṁ modalaindi.}
\end{quote}

\begin{quote}
When the tickets got mixed up, confusion broke out at the station.
\end{quote}

\te{గందరగోళంగా} \te{ఉంది}, it is confusing; \te{గందరగోళం} \te{చేయు}, to make a mess of things.

\begin{cousinweb}
\textbf{\te{గందరగోళం}} (\emph{gandaragōḷaṁ}) is an everyday Telugu word for confusion and uproar whose origin is not settled.
\end{cousinweb}

\begin{grammarlens}[title={What you've built}]
One more word for getting around.
\end{grammarlens}

\subsection*{Your turn}
\begin{itemize}
\raggedright
  \item \emph{Say it:} \emph{gandaragōḷaṁ}
  \item \emph{Say it:} \emph{gandaragōḷaṁ}, once more
  \item \emph{Say it:} \emph{ṭikkeṭlu tārumāru kāvaḍantō sṭēṣanlō gandaragōḷaṁ modalaindi.} — When the tickets got mixed up, confusion broke out at the station.
  \item \emph{Recall:} say the Telugu for to estimate, then the Telugu for a false step, then say \emph{gandaragōḷaṁ} again
\end{itemize}

\begin{tcolorbox}[breakable,colback=teal!4,colframe=teal!35!black,title={Before you move on}]
What does \te{గందరగోళం} mean? (\textquotedblleft{}Confusion; a muddle\textquotedblright{}.) Say it once more.
\end{tcolorbox}
